\documentclass{beamer}

\usepackage{luatexja}
\usepackage{luatexja-fontspec}

% ---- ブロックの設定（block / alertblock / exampleblock）----
\definecolor{alertcol}{RGB}{200,80,40}            % alertblock の色（オレンジ系）
\definecolor{examplecol}{RGB}{30,120,60}          % exampleblock の色（緑系）
\setbeamertemplate{blocks}[rounded][shadow=false] % 角丸・影なし（3種類すべてに適用）

% block（通常）：青
\setbeamercolor{block title}{fg=white,bg=main}
\setbeamercolor{block body}{fg=black,bg=main!6}

% alertblock（注意）：オレンジ
\setbeamercolor{block title alerted}{fg=white,bg=alertcol}
\setbeamercolor{block body alerted}{fg=black,bg=alertcol!6}

% exampleblock（例）：緑
\setbeamercolor{block title example}{fg=white,bg=examplecol}
\setbeamercolor{block body example}{fg=black,bg=examplecol!6}

% タイトルを太字に（任意）
\setbeamerfont{block title}{series=\bfseries}
\setbeamerfont{block title alerted}{series=\bfseries}
\setbeamerfont{block title example}{series=\bfseries}
%====================================================
% --------------------------------------------------
\setbeamertemplate{itemize items}[circle]   % 平たい丸
% \setbeamertemplate{itemize items}[ball]     % 立体的な丸
% \setbeamertemplate{itemize items}[square]   % 四角
% \setbeamertemplate{itemize items}[triangle] % 三角（今の状態）
%==================================================
\definecolor{main}{RGB}{30,70,140}                  % バナーで定義済みなら不要
\setbeamercolor{block title}{fg=white,bg=main}      % タイトル部：白文字・青背景
\setbeamercolor{block body}{bg=main!6}              % 本体：main を 6% に薄めた色
\setbeamertemplate{blocks}[rounded][shadow=false]   % 角を丸める・影なし
\setbeamertemplate{itemize items}[circle]           % 箇条書きを丸に

\newcommand{\argmax}{\mathop{\rm arg~max}\limits}
\newcommand{\argmin}{\mathop{\rm arg~min}\limits}
%===================================================

\title{The Principles of Diffusion Models 1-1}
\author{渡邉 有志}
\date{\today}

\begin{document}

\begin{frame}
    \titlepage
\end{frame}

\begin{frame}{目次}
    
    \begin{itemize}
        \item 
    \end{itemize}
\end{frame}

%=======================================================
\begin{frame}{DGM(深層生成モデル)}
\begin{block}{DGM(Deep Generative Models)}
高次元データに対する確率分布を学習し、データセットに類似したサンプルの生成が可能なモデル\\
\end{block}

\vspace{-1em}           

\begin{columns}[T]
\begin{column}{0.48\textwidth}
\begin{exampleblock}{\small 目標1 Realistic Generation}
\small 実際のデータと区別できないような\textbf{新しいサンプルを生成}
\end{exampleblock}
\end{column}

\begin{column}{0.48\textwidth}
\begin{exampleblock}{\small 目標2 Controllable Generation}
\small 生成過程を解釈可能な形での制御を可能に
\end{exampleblock}
\end{column}
\end{columns}

\vspace{1em}     

\small モデルが表す確率分布:\boldmath{$p_{\phi}$}  実際のデータが従う確率分布\footnote{data distribution, true distributionとも呼び、一般に$p_{data}$自体は未知である}:$p_{data}$\unboldmath\\

\begin{itemize}
\item 有限個のデータセットに対し\\
$p_{\phi}$と$p_{data}$間の損失関数を最小化→${\phi}$を最適化
\end{itemize}

\begin{exampleblock}{}
⇒つまり$p_{\phi}(x)\approx p_{data}(x)$が理想的\\
しかし$p_{data}(x)$が不明なので手元にある$p_{data}$から得られた\\サンプル$x^{(1)},x^{(2)},...x^{(N)}$から$\phi$を学習 ⇒ 分布を推定
\end{exampleblock}
\end{frame}
%=======================================================
\begin{frame}{1.1.1 Mathematical Setup -Goal of DGM}
\begin{exampleblock}{}
    前提：未知の分布$p_{data}$からサンプリングした有限個のサンプル\[
\left\{\mathbf{x}^{(i)}\right\}_{i=1}^{N}が手元にある
\]
\end{exampleblock}
\small $p_{data}$を表す式、新しいサンプルを生成する方法もわからない\\
⇒そこでパラメータ$\phi$を持つ深層ネットで\textbf{モデル分布}$p_{\phi}(x)$を構成\footnote{pは文脈により($p_{data}$)確率分布または($p_{data}(x)$)確率密度関数}
\begin{block}{DGMの目標}
    \centering$\phi^{*}$を学習して$p_{\phi^{*}}(\boldsymbol{x})\approx p_{data}(\boldsymbol{x})$\\
    $p_{\phi}$をGenerative model（生成モデル）という
\end{block}
⇒未知の$p_{data}$から得られた有限データを使って$p_{\phi^{*}}\approx p_{data}$となる\\
モデルを学習し、そこから新しいデータを生成する

\end{frame}
%=======================================================
\begin{frame}{1-1-1 Training of DGM}
\begin{block}{学習}
モデル族$\{p_{\phi}\}$の中で不一致$\mathcal{D}$を最小化するパラメータを求める\\
\centering $\phi^{*}\in \argmin_{\phi}\mathcal{D}(p_{data},p_{\phi})$
\end{block}


\end{frame}
%=======================================================
\begin{frame}{Forward KL divergence}
\begin{block}{義:KL divergence}
    $D_{KL}(p_{data}\mid\mid p_{\phi}):=\int p_{data}({\boldsymbol{x}})log\frac{p_{data}(\boldsymbol{x})}{p_{\phi}(\boldsymbol{x})}dx=\boldsymbol{\mathbb{E}_{\boldsymbol{x}\sim p_{data}}[logp_{data}(\boldsymbol{x})-logp_{\phi}(\boldsymbol{x})]}$
\end{block}

\begin{itemize}
    \item 非対称:$D_{KL}(p_{data}\mid p_{\phi} )\neq D_{KL}(p_{\phi}\mid\mid p_{data})$
    \item 
    
\end{itemize}
\end{frame}
%=======================================================
\begin{frame}

\end{frame}
%=======================================================
%=======================================================
%=======================================================
\end{document}
